\begin{abstract}
\noindent
Scanned PDFs and image-based documents are ubiquitous, yet many of them carry no
machine-readable text, so the browser's familiar ``find on page''
($\mathrm{Ctrl}{+}\mathrm{F}$) returns nothing. This report presents \proj{}, a
Chrome extension (Manifest~V3) that closes this gap entirely on the user's
device \emph{without replacing anything the user sees}. Chrome's own PDF viewer
and its own find bar remain the only interface. When the user presses a shortcut
on a scanned PDF, a hidden offscreen document downloads the file, detects which
pages paint images or lack text, recognizes those pages with Tesseract.js, and
writes every recognized word back into a \emph{copy} of the PDF as invisible text
(PDF text rendering mode~3), sized so that the native find highlight covers the
visible word. The tab is then switched to that copy, which the extension's service
worker serves from Cache Storage, so Chrome opens it in its native viewer with
the same tab title. The original file is never modified, and the copy renders
pixel-identically.

The report follows the project from a custom PDF.js viewer with an injected DOM
text layer to this ``searchable-scan sandwich'' design, and explains why the
change was needed: no DOM viewer can reproduce Chrome's native PDF experience.
It describes the coordinate mapping from OCR pixels into PDF user space, the OCR
preprocessing (including a gain-capped contrast stretch that keeps sparse pages
from collapsing into noise), content-addressed caching in IndexedDB, and
resource management for a long-lived background pipeline. It closes with an
end-to-end test suite that drives the packaged extension in real Chrome. That
suite exposed a defect that unit-level tests had missed: Chrome refuses to
navigate an ordinary web tab to an extension-owned \code{blob:} URL, so the
first version never actually switched tabs.
\end{abstract}
